%=======================================================================
%  Algorithms: An Analytical Note
%  Discrete Mathematics (MSc IT) -- Assignment
%  Student : Solanki Ruchit M.   ID: 202612085
%  Course  : MSc IT
%  Instructor : Prof. Gopinath Panda
%  Subject : Discrete Mathematics
%  Reference : Kenneth H. Rosen, "Discrete Mathematics and Its
%              Applications", 8th edition, Chapter 3 (Algorithms).
%
%  Compile with :  pdflatex  (or xelatex / lualatex)
%=======================================================================
\documentclass[11pt,a4paper]{article}

\usepackage[a4paper,top=0.85in,bottom=0.85in,left=0.95in,right=0.95in]{geometry}
\usepackage{amsmath,amssymb,amsthm}
\usepackage{graphicx}
\usepackage{booktabs}
\usepackage{array}
\usepackage{enumitem}
\usepackage{xcolor}
\usepackage{titlesec}
\usepackage{fancyhdr}
\usepackage{caption}
\usepackage[hidelinks]{hyperref}

%------------------------------------------------------------- styling
\definecolor{navy}{HTML}{12355B}
\definecolor{accent}{HTML}{1F6FEB}
\definecolor{teal}{HTML}{0E9594}
\definecolor{softgray}{HTML}{F4F6F9}

\titleformat{\section}{\normalfont\large\bfseries\color{navy}}{\thesection.}{0.6em}{}
\titleformat{\subsection}{\normalfont\normalsize\bfseries\color{accent}}{\thesubsection}{0.6em}{}
\titlespacing*{\section}{0pt}{0.8em}{0.35em}
\titlespacing*{\subsection}{0pt}{0.6em}{0.25em}
\setlength{\abovedisplayskip}{4pt}
\setlength{\belowdisplayskip}{4pt}

\pagestyle{fancy}
\fancyhf{}
\fancyhead[L]{\footnotesize Solanki Ruchit M. \textbar\ 202612085}
\fancyhead[R]{\footnotesize Discrete Mathematics \textbar\ Algorithms}
\fancyfoot[C]{\footnotesize\thepage}
\renewcommand{\headrulewidth}{0.4pt}

\captionsetup{font=small,labelfont={bf,color=navy},skip=3pt}
\setlength{\textfloatsep}{8pt plus 2pt minus 2pt}
\setlength{\intextsep}{8pt plus 2pt minus 2pt}

%------------------------------------------------------------- theorems
\theoremstyle{definition}
\newtheorem{definition}{Definition}[section]
\newtheorem{example}{Example}[section]
\theoremstyle{plain}
\newtheorem{theorem}{Theorem}[section]
\newtheorem{corollary}{Corollary}[section]

\newcommand{\bigO}{\mathcal{O}}
\newcommand{\bigOmega}{\Omega}
\newcommand{\bigTheta}{\Theta}

\setlist[itemize]{leftmargin=1.3em,itemsep=2pt,topsep=3pt}
\setlist[enumerate]{leftmargin=1.5em,itemsep=2pt,topsep=3pt}

\begin{document}

%------------------------------------------------------------- title
\begin{center}
{\LARGE\bfseries\color{navy} Algorithms: An Analytical Note}\\[6pt]
{\large Discrete Mathematics --- Assignment}\\[10pt]
{\small
\begin{tabular}{@{}ll@{}}
\textbf{Student:} & Solanki Ruchit M. \quad (ID: 202612085)\\
\textbf{Course:} & MSc IT\\
\textbf{Subject:} & Discrete Mathematics\\
\textbf{Instructor:} & Prof. Gopinath Panda\\
\textbf{Reference:} & K.~H. Rosen, \emph{Discrete Mathematics and Its}\\
& \emph{Applications}, 8th ed., Chapter~3 (Algorithms)\\
\end{tabular}}
\end{center}
\vspace{-4pt}
\hrule
\vspace{6pt}

\begin{abstract}
\noindent
This note presents a complete, self--contained treatment of \emph{algorithms} as
developed in Chapter~3 of Rosen's \emph{Discrete Mathematics and Its Applications}.
It answers four questions in order: \textbf{(i)}~what exactly an algorithm is and
which properties make a procedure an algorithm; \textbf{(ii)}~the principal
\emph{types} of algorithms (searching, sorting, greedy and divide--and--conquer);
\textbf{(iii)}~how the \emph{growth of functions} is measured formally with the
$\bigO$, $\bigOmega$ and $\bigTheta$ notations; and \textbf{(iv)}~how
\emph{complexity} is analysed in the worst, average and best cases, together with
the notion of tractability. Every definition is stated formally and illustrated
with a small diagram, and each claim about running time is justified
analytically. The companion artefacts of this submission are the LaTeX
pseudocode file, the executed Jupyter notebook and the interactive HTML
dashboard.
\end{abstract}

%=======================================================================
\section{Introduction: What is an Algorithm?}
%=======================================================================

The word \emph{algorithm} derives from the ninth--century Persian mathematician
Muhammad ibn Musa \emph{al--Khwarizmi}; the Latinised form \emph{Algoritmi}
became \emph{algorithm}. Intuitively an algorithm is a recipe, but discrete
mathematics demands a precise definition.

\begin{definition}[Algorithm, Rosen \S3.1]
An \textbf{algorithm} is a finite sequence of precise instructions for performing
a computation or for solving a problem.
\end{definition}

\begin{definition}[The seven properties]
A procedure is an algorithm only if it satisfies all of the following:
\begin{enumerate}[label=\textbf{(\arabic*)}]
\item \textbf{Input} --- it receives a finite set of values from a specified set;
\item \textbf{Output} --- it produces a solution from a specified set;
\item \textbf{Definiteness} --- each step is precisely and unambiguously stated;
\item \textbf{Correctness} --- it produces the correct output for every legal input;
\item \textbf{Finiteness} --- it produces the output after a \emph{finite} number of steps;
\item \textbf{Effectiveness} --- each step is basic enough to be carried out exactly;
\item \textbf{Generality} --- it solves a whole \emph{class} of problems, not one instance.
\end{enumerate}
\end{definition}

Figure~\ref{fig:box} summarises this view of an algorithm as a transducer of
inputs into outputs.

\begin{figure}[htbp]
\centering
\includegraphics[width=0.95\textwidth]{figures/fig1_algorithm_box.pdf}
\caption{An algorithm transforms an input instance into an output solution, and
must satisfy the seven properties of Definition~1.2.}
\label{fig:box}
\end{figure}

\subsection{Pseudocode conventions}
Algorithms in Chapter~3 are written in \emph{pseudocode}, an informal but
unambiguous mixture of English and programming syntax. Its conventions are:
assignment \texttt{:=}; block structure by indentation; control flow with
\texttt{if--then--else}, \texttt{while}, \texttt{for}; procedures returning a
value via \texttt{return}; and \texttt{\{ \}} for a comment. Because pseudocode
hides the machine, it exposes the \emph{logic} and therefore the \emph{cost} of a
procedure --- the central concern of \S3 and \S4 below.

\subsection{The limits of algorithms}
Not every problem has an algorithmic solution. The \textbf{halting problem} asks
for a procedure that decides, for an arbitrary program and input, whether the
program terminates. Turing proved that \emph{no such algorithm exists}: the
problem is \emph{undecidable}. This is the canonical reason the definition above
insists on \emph{definiteness} and \emph{finiteness} --- an algorithm is not merely
``a way of computing'', but a guaranteed termination with a specified answer.

%=======================================================================
\section{Types of Algorithms}
%=======================================================================

Chapter~3 presents several \emph{paradigms}. A paradigm is a strategy for
designing an algorithm, and the paradigm usually determines the growth rate of
the resulting procedure.

\subsection{Searching algorithms}
\textbf{Linear search} examines the list elements one after another until the
target is found or the list is exhausted. \textbf{Binary search} works only on a
list that is \emph{already sorted}: it compares the target with the middle
element and discards half of the remaining list at every step. The two behave
completely differently, as Figure~\ref{fig:search} shows.

\begin{figure}[htbp]
\centering
\includegraphics[width=0.95\textwidth]{figures/fig2_search_comparison.pdf}
\caption{Searching for 61 in a sorted list of 12 elements. (a) Linear search
performs 9 comparisons; (b) binary search performs only 3, because each step
halves the interval.}
\label{fig:search}
\end{figure}

\subsection{Sorting algorithms}
Sorting rearranges a list into increasing order. Three elementary
$\bigO(n^2)$ sorts are treated in Chapter~3:
\begin{itemize}
\item \textbf{Bubble sort} repeatedly compares adjacent pairs and swaps them when
out of order; each pass moves the largest remaining value to the end
(Figure~\ref{fig:bubble}).
\item \textbf{Insertion sort} maintains a sorted prefix and inserts the next
element into its correct place.
\item \textbf{Selection sort} repeatedly selects the smallest remaining element
and appends it to the sorted prefix.
\end{itemize}
A faster divide--and--conquer sort, \textbf{merge sort}, splits the list in half,
sorts each half recursively and merges them in linear time, giving
$\bigO(n\log n)$.

\begin{figure}[htbp]
\centering
\includegraphics[width=0.95\textwidth]{figures/fig3_bubble_sort.pdf}
\caption{Bubble sort applied to $[29,10,14,37,13]$. Shaded cells are already in
final position. The number of comparisons is exactly $n(n-1)/2$, hence
$\bigO(n^2)$.}
\label{fig:bubble}
\end{figure}

\subsection{Greedy algorithms}
A \textbf{greedy algorithm} builds a solution one step at a time, always choosing
the option that looks best \emph{right now}. The classical illustration is
\emph{change--making}: to give change of $n$ cents using
$25,10,5,1$--cent coins, repeatedly take the largest coin not exceeding the
remaining amount. Greedy is fast (one pass over the denominations) but it is
\emph{not always optimal} for arbitrary coin systems --- the U.S.\ system is one
of the lucky ones.

\subsection{Recursive and divide--and--conquer algorithms}
A \textbf{recursive} algorithm calls itself on a smaller instance. When the
recursion splits the input into independent parts, the paradigm is called
\emph{divide--and--conquer}; its cost is described by a \emph{recurrence
relation}. Merge sort satisfies $T(n)=2T(n/2)+n$, whose solution is
$\bigTheta(n\log n)$.

\begin{table}[htbp]
\centering
\small
\caption{Summary of the algorithm types discussed in Chapter~3.}
\label{tab:types}
\begin{tabular}{@{}p{3.4cm}p{4.4cm}p{4.2cm}c@{}}
\toprule
\textbf{Type} & \textbf{Idea} & \textbf{Typical example} & \textbf{Growth}\\
\midrule
Linear search & scan every element & find a key in an unsorted list & $\bigO(n)$\\
Binary search & halve a sorted interval & find a key in a sorted list & $\bigO(\log n)$\\
Bubble / insertion / selection sort & repeated swaps or insertions & sort a small list & $\bigO(n^2)$\\
Merge sort & divide, conquer, merge & sort a large list & $\bigO(n\log n)$\\
Greedy & make the locally best choice & coin change, scheduling & $\bigO(n)$--$\bigO(n\log n)$\\
Recursive & reduce to a smaller instance & Towers of Hanoi, factorials & $2^n$, $n!$, \dots\\
\bottomrule
\end{tabular}
\end{table}

%=======================================================================
\section{The Growth of Functions}
%=======================================================================

To compare algorithms we must compare the \emph{growth} of the functions that
count their operations as the input size $n$ grows. Chapter~3 formalises this
intuition with the ``big--$\bigO$'' family of notations.

\begin{definition}[Big--$\bigO$]
Let $f,g:\mathbb{Z}\to\mathbb{R}$. We say $f(x)$ is $\bigO(g(x))$ if there exist
\emph{witnesses} $C>0$ and $k>0$ such that
\[
|f(x)|\le C\,|g(x)| \qquad \text{for all } x>k .
\]
\end{definition}

\begin{definition}[Big--$\bigOmega$ and Big--$\bigTheta$]
$f(x)$ is $\bigOmega(g(x))$ if $|f(x)|\ge C\,|g(x)|$ for all $x>k$ (a lower
bound). We write $f(x)$ is $\bigTheta(g(x))$ if $f(x)$ is both
$\bigO(g(x))$ and $\bigOmega(g(x))$, i.e.\ $f$ and $g$ grow at the same rate.
\end{definition}

The witnesses $C$ and $k$ are the whole content of the definition: $\bigO$ makes
a claim about the \emph{tail} of the function, not about small values.

\begin{theorem}[Sum and product rules]
If $f_1=\bigO(g_1)$ and $f_2=\bigO(g_2)$, then
\[
(f_1+f_2)(x)=\bigO\!\big(\max(|g_1(x)|,|g_2(x)|)\big),
\qquad
(f_1f_2)(x)=\bigO\!\big(g_1(x)g_2(x)\big).
\]
\end{theorem}

\begin{theorem}[Polynomials]
If $a_0,a_1,\dots,a_k$ are real with $a_k\neq0$, then
\[
a_k n^k+a_{k-1}n^{k-1}+\dots+a_0 \;=\; \bigO(n^k),
\]
so the leading term alone determines the growth of a polynomial.
\end{theorem}

Figure~\ref{fig:growth} plots the standard functions. On a linear scale the
exponential $2^n$ destroys every polynomial; on a logarithmic scale the powers
$n$, $n^2$, $n^3$ become straight lines, which is exactly why logarithmic plots
are used to compare growth rates.

\begin{figure}[htbp]
\centering
\includegraphics[width=0.95\textwidth]{figures/fig4_growth.pdf}
\caption{Growth of the standard complexity functions on a linear axis (left) and
a logarithmic axis (right).}
\label{fig:growth}
\end{figure}

\begin{figure}[htbp]
\centering
\includegraphics[width=0.92\textwidth]{figures/fig5_bigo_hierarchy.pdf}
\caption{Inclusion hierarchy of the common complexity classes. Each class
contains all the classes drawn inside it.}
\label{fig:hierarchy}
\end{figure}

\begin{table}[htbp]
\centering
\small
\caption{Numerical values of the standard growth functions.}
\label{tab:values}
\begin{tabular}{@{}rrrrrr@{}}
\toprule
$n$ & $\log_2 n$ & $n$ & $n\log_2 n$ & $n^2$ & $2^n$\\
\midrule
$10$      & $3.3$  & $10$    & $33$      & $10^{2}$ & $1.02\times10^{3}$\\
$10^{2}$  & $6.6$  & $10^{2}$& $6.6\times10^{2}$ & $10^{4}$ & $1.27\times10^{30}$\\
$10^{3}$  & $10$   & $10^{3}$& $10^{4}$  & $10^{6}$ & $1.07\times10^{301}$\\
$10^{4}$  & $13.3$ & $10^{4}$& $1.3\times10^{5}$ & $10^{8}$ & astronomically large\\
\bottomrule
\end{tabular}
\end{table}

%=======================================================================
\section{Complexity of Algorithms}
%=======================================================================

\subsection{Time complexity}
The \textbf{time complexity} of an algorithm is the number of basic operations it
performs, expressed as a function of the input size $n$ using the notation of
\S3. Three cases are distinguished:
\begin{itemize}
\item \textbf{Worst--case} complexity $W(n)$ --- the maximum over all inputs of
size $n$; the most useful guarantee.
\item \textbf{Average--case} complexity $A(n)$ --- the expected cost assuming a
probability distribution on inputs.
\item \textbf{Best--case} complexity $B(n)$ --- the minimum; usually of little
practical value.
\end{itemize}

\subsection{Worked complexity analysis}
\noindent\textbf{Linear search.} In the worst case the target is absent and every
element is tested; the cost is $\bigTheta(n)$.\\
\textbf{Binary search.} After $i$ steps the interval has size $n/2^{\,i-1}$; it
becomes empty when $n/2^{\,i-1}<1$, i.e.\ $i>\log_2 n+1$. Hence
$W(n)=\bigTheta(\log n)$ --- the enormous saving visible in Figure~\ref{fig:search}.\\
\textbf{Bubble sort.} The inner loop performs comparisons
\[
\sum_{i=1}^{n-1}(n-i)=\frac{n(n-1)}{2}=\bigTheta(n^2),
\]
so bubble sort is quadratic in every case.

\subsection{Space complexity and tractability}
\textbf{Space complexity} counts memory cells rather than operations. A problem is
called \emph{tractable} if it can be solved by an algorithm with polynomial
worst--case time complexity; problems requiring exponential time are deemed
\emph{intractable}, because for even modest $n$ the count explodes
(Figure~\ref{fig:bars} and Table~\ref{tab:values}). The question of whether every
problem whose solution can be \emph{checked} quickly can also be \emph{solved}
quickly is the famous $\mathrm{P}$ versus $\mathrm{NP}$ problem.

\begin{figure}[htbp]
\centering
\includegraphics[width=0.95\textwidth]{figures/fig6_complexity_bars.pdf}
\caption{Basic--operation counts at growing input sizes. The polynomial methods
stay feasible while the number of operations for exponential methods becomes
astronomical.}
\label{fig:bars}
\end{figure}

%=======================================================================
\section{Worked Analytical Examples}
%=======================================================================

\begin{example}[Finding the witnesses]
Show that $f(n)=3n^2+2n+5$ is $\bigO(n^2)$.
\end{example}
\begin{proof}
For every $n\ge1$ we have $2n\le 2n^2$ and $5\le 5n^2$. Therefore
\[
f(n)=3n^2+2n+5 \le 3n^2+2n^2+5n^2 = 10n^2 .
\]
Taking $C=10$ and $k=1$ gives $f(n)\le 10n^2$ for all $n>1$, so
$f(n)=\bigO(n^2)$. Moreover $f(n)\ge 3n^2$ for $n\ge1$, so
$f(n)=\bigOmega(n^2)$; combining the two bounds, $f(n)=\bigTheta(n^2)$.
\end{proof}

\begin{example}[Separating two growth rates]
Show that $n^2$ is \emph{not} $\bigO(n\log n)$.
\end{example}
\begin{proof}
Suppose, for contradiction, that $n^2\le C\,n\log_2 n$ for all $n>k$. Dividing by
$n\log_2 n$ gives
\[
\frac{n}{\log_2 n}\le C \qquad\text{for all } n>k .
\]
But $n/\log_2 n\to\infty$ as $n\to\infty$; choosing $n>2^{\,C+1}$ makes the
quotient exceed $C$, a contradiction. Hence $n^2\neq\bigO(n\log n)$, which is
precisely why merge sort beats bubble sort on large inputs. \qedhere
\end{proof}

\begin{example}[Binary search recurrence]
Let $W(n)$ be the worst--case number of comparisons of binary search. Then
\[
W(n)=W(\lfloor n/2\rfloor)+1,\qquad W(1)=1 .
\]
Unrolling the recurrence, $W(n)=1+\lfloor\log_2 n\rfloor$, i.e.\ $W(n)=\bigTheta(\log n)$.
\end{example}

%=======================================================================
\section{Conclusion}
%=======================================================================

An algorithm is a finite, definite, effective and general procedure whose
correctness is guaranteed for every legal input. Its power is measured not by
how clever it looks but by how its operation count \emph{grows}: the
$\bigO$--$\bigOmega$--$\bigTheta$ calculus of \S3 turns ``this algorithm feels
faster'' into a theorem, and the worst--case analysis of \S4 decides whether a
problem is tractable. The four paradigms of \S2 illustrate the whole range: a
single change of strategy --- scanning versus halving, or quadratic versus
divide--and--conquer --- moves a procedure from $\bigO(n^2)$ to
$\bigO(n\log n)$ or from $\bigO(n)$ to $\bigO(\log n)$, which for large inputs is
the difference between feasible and impossible.

%=======================================================================
\section*{References}
\addcontentsline{toc}{section}{References}
%=======================================================================
\begin{enumerate}[label={[\arabic*]},leftmargin=2.2em]
\item K.~H. Rosen, \emph{Discrete Mathematics and Its Applications},
      8th ed., McGraw--Hill, 2019. \textbf{Chapter~3: Algorithms} (pp.~191--242).
\item T.~H. Cormen, C.~E. Leiserson, R.~L. Rivest and C.~Stein,
      \emph{Introduction to Algorithms}, 3rd ed., MIT Press, 2009.
\item D.~E. Knuth, \emph{The Art of Computer Programming, Vol.~1:
      Fundamental Algorithms}, 3rd ed., Addison--Wesley, 1997.
\item R.~Sedgewick and K.~Wayne, \emph{Algorithms}, 4th ed., Addison--Wesley, 2011.
\item J.~Kleinberg and E.~Tardos, \emph{Algorithm Design},
      Pearson/Addison--Wesley, 2006.
\item A.~M. Turing, ``On computable numbers, with an application to the
      Entscheidungsproblem'', \emph{Proc.\ London Math.\ Soc.}, ser.~2,
      vol.~42, pp.~230--265, 1936.
\end{enumerate}

\end{document}
